\documentclass[11pt]{article}

\usepackage[margin=1in]{geometry}
\usepackage{amsmath,amssymb,amsthm}
\usepackage{bm}
\usepackage{booktabs}
\usepackage{setspace}
\usepackage{enumitem}
\usepackage{tcolorbox}

\onehalfspacing

\newcommand{\Var}{\operatorname{Var}}
\newcommand{\Cov}{\operatorname{Cov}}

\title{The Bivariate Normal CEF\\[4pt]\large Why the CEF is Linear Under Joint Normality}
\author{PLSC 30700: Estimation I}
\date{}

\begin{document}

\maketitle

\section{Setup}

Suppose two random variables are jointly bivariate normal:
\[
\begin{pmatrix} X \\ Y \end{pmatrix}
\sim N\!\left(
\begin{pmatrix} \mu_X \\ \mu_Y \end{pmatrix},\
\begin{pmatrix} \sigma_X^2 & \rho\,\sigma_X\sigma_Y \\ \rho\,\sigma_X\sigma_Y & \sigma_Y^2 \end{pmatrix}
\right)
\]
where $\rho = \Cov(X,Y)/(\sigma_X\sigma_Y)$ is the correlation.  Our goal is to derive the conditional distribution $f_{Y|X}(y\mid x)$ and show that $\mathbb{E}[Y\mid X=x]$ is a \emph{linear} function of $x$.  This is one of the classical motivations for linear regression.

\medskip
Throughout, we write densities as functions of realized values $(x,y)$.  Define the standardized values
\[
x^* = \frac{x - \mu_X}{\sigma_X}, \qquad y^* = \frac{y - \mu_Y}{\sigma_Y}.
\]
(These are numbers, not random variables.  The corresponding random variables would be $X^*=(X-\mu_X)/\sigma_X$ and $Y^*=(Y-\mu_Y)/\sigma_Y$, but inside density expressions we always evaluate at a specific point.)

\section{The three densities we need}

\subsection{Marginal density of $X$}

Since $X\sim N(\mu_X, \sigma_X^2)$:
\begin{equation}\label{eq:marginal}
f_X(x) = \frac{1}{\sqrt{2\pi}\,\sigma_X}\,\exp\!\left(-\frac{1}{2}(x^*)^2\right)
\end{equation}

\subsection{Joint density of $(X,Y)$}

The bivariate normal density is:
\begin{equation}\label{eq:joint}
f_{X,Y}(x,y)
= \frac{1}{2\pi\,\sigma_X\sigma_Y\sqrt{1-\rho^2}}\,
\exp\!\left(
-\frac{1}{2(1-\rho^2)}
\Big[(x^*)^2 - 2\rho\, x^* y^* + (y^*)^2\Big]
\right)
\end{equation}

\subsection{Conditional density: the ratio}

By definition:
\[
f_{Y|X}(y\mid x) = \frac{f_{X,Y}(x,y)}{f_X(x)}
\]
We will compute this ratio and simplify.

\section{Step 1: Form the ratio}

Dividing \eqref{eq:joint} by \eqref{eq:marginal}:
\[
f_{Y|X}(y\mid x)
= \frac{
\dfrac{1}{2\pi\,\sigma_X\sigma_Y\sqrt{1-\rho^2}}\,
\exp\!\left(
-\dfrac{(x^*)^2 - 2\rho\, x^* y^* + (y^*)^2}{2(1-\rho^2)}
\right)
}{
\dfrac{1}{\sqrt{2\pi}\,\sigma_X}\,
\exp\!\left(-\dfrac{(x^*)^2}{2}\right)
}
\]

\medskip\noindent
\textbf{Simplify the constant out front.}  The ratio of the leading fractions is:
\[
\frac{1/(2\pi\,\sigma_X\sigma_Y\sqrt{1-\rho^2})}{1/(\sqrt{2\pi}\,\sigma_X)}
= \frac{\sqrt{2\pi}\,\sigma_X}{2\pi\,\sigma_X\sigma_Y\sqrt{1-\rho^2}}
= \frac{1}{\sqrt{2\pi}\,\sigma_Y\sqrt{1-\rho^2}}
\]

\medskip\noindent
\textbf{Simplify the exponential.}  Using $e^a / e^b = e^{a-b}$, the exponent becomes:
\[
-\frac{(x^*)^2 - 2\rho\, x^* y^* + (y^*)^2}{2(1-\rho^2)}
- \left(-\frac{(x^*)^2}{2}\right)
\]
\[
= -\frac{(x^*)^2 - 2\rho\, x^* y^* + (y^*)^2}{2(1-\rho^2)}
+ \frac{(x^*)^2}{2}
\]

\medskip\noindent
Put the second term over the common denominator $2(1-\rho^2)$:
\[
= -\frac{(x^*)^2 - 2\rho\, x^* y^* + (y^*)^2}{2(1-\rho^2)}
+ \frac{(1-\rho^2)(x^*)^2}{2(1-\rho^2)}
\]
\[
= \frac{-(x^*)^2 + 2\rho\, x^* y^* - (y^*)^2 + (x^*)^2 - \rho^2(x^*)^2}{2(1-\rho^2)}
\]

\medskip\noindent
The $(x^*)^2$ terms cancel:
\[
= \frac{2\rho\, x^* y^* - (y^*)^2 - \rho^2(x^*)^2}{2(1-\rho^2)}
\]

\medskip\noindent
Factor out a negative sign:
\[
= -\frac{(y^*)^2 - 2\rho\, x^* y^* + \rho^2(x^*)^2}{2(1-\rho^2)}
\]

\section{Step 2: Recognize the perfect square}

The numerator in the exponent is:
\[
(y^*)^2 - 2\rho\, x^* y^* + \rho^2 (x^*)^2 = (y^* - \rho\, x^*)^2
\]
You can verify this by expanding the right side:
\[
(y^* - \rho\, x^*)^2 = (y^*)^2 - 2\rho\, x^* y^* + \rho^2(x^*)^2 \quad\checkmark
\]

So the exponent simplifies to:
\[
-\frac{(y^* - \rho\, x^*)^2}{2(1-\rho^2)}
\]

\medskip\noindent
Putting it all together:
\begin{equation}\label{eq:conditional-star}
f_{Y|X}(y\mid x)
= \frac{1}{\sqrt{2\pi}\,\sigma_Y\sqrt{1-\rho^2}}\,
\exp\!\left(
-\frac{(y^* - \rho\, x^*)^2}{2(1-\rho^2)}
\right)
\end{equation}

\section{Step 3: Rewrite in terms of $y$ and $x$}

Substitute $y^* = (y-\mu_Y)/\sigma_Y$ and $x^* = (x-\mu_X)/\sigma_X$ back in:
\[
y^* - \rho\, x^*
= \frac{y - \mu_Y}{\sigma_Y} - \rho\,\frac{x - \mu_X}{\sigma_X}
= \frac{1}{\sigma_Y}\left[y - \mu_Y - \rho\,\frac{\sigma_Y}{\sigma_X}(x-\mu_X)\right]
\]

\medskip\noindent
So the exponent in \eqref{eq:conditional-star} becomes:
\[
-\frac{1}{2(1-\rho^2)}\left(\frac{y - \mu_Y - \rho\frac{\sigma_Y}{\sigma_X}(x-\mu_X)}{\sigma_Y}\right)^{\!2}
= -\frac{1}{2}\left(\frac{y - \mu_Y - \rho\frac{\sigma_Y}{\sigma_X}(x-\mu_X)}{\sigma_Y\sqrt{1-\rho^2}}\right)^{\!2}
\]

\medskip\noindent
The full conditional density is therefore:
\begin{equation}\label{eq:conditional-full}
f_{Y|X}(y\mid x)
= \frac{1}{\sqrt{2\pi}\;\sigma_Y\sqrt{1-\rho^2}}\;
\exp\!\left(
-\frac{1}{2}\left(
\frac{y - \Big[\mu_Y + \rho\frac{\sigma_Y}{\sigma_X}(x - \mu_X)\Big]}{\sigma_Y\sqrt{1-\rho^2}}
\right)^{\!2}
\right)
\end{equation}

\section{Step 4: Read off the conditional distribution}

Equation~\eqref{eq:conditional-full} is the density of a normal distribution.  Recall that $N(\mu, \tau^2)$ has density
\[
\frac{1}{\sqrt{2\pi}\,\tau}\exp\!\left(-\frac{(y-\mu)^2}{2\tau^2}\right).
\]
Matching terms:

\begin{tcolorbox}[colback=blue!5, colframe=blue!50]
\textbf{Conditional distribution:}\quad $Y\mid X=x \;\sim\; N\!\Big(\mu_{Y|X},\;\sigma^2_{Y|X}\Big)$

\medskip
\begin{itemize}[leftmargin=2em]
\item \textbf{Conditional mean:}\quad
$\displaystyle\mathbb{E}[Y\mid X=x] = \mu_Y + \rho\,\frac{\sigma_Y}{\sigma_X}\,(x - \mu_X)$

\medskip
\item \textbf{Conditional variance:}\quad
$\displaystyle\Var(Y\mid X=x) = \sigma_Y^2(1 - \rho^2)$
\end{itemize}
\end{tcolorbox}

\section{Step 5: Connect to regression}

\subsection{The CEF is linear}

The conditional expectation function is:
\[
\mathbb{E}[Y \mid X = x]
= \underbrace{\mu_Y - \rho\,\frac{\sigma_Y}{\sigma_X}\,\mu_X}_{\text{intercept}} \;+\; \underbrace{\rho\,\frac{\sigma_Y}{\sigma_X}}_{\text{slope}}\cdot x
\]
This is a linear function of $x$.  Under joint normality, the CEF is \emph{exactly} linear---no approximation involved.

\subsection{The slope is the regression coefficient}

Recall that the population regression coefficient from the BLP is:
\[
\beta = \frac{\Cov(Y,X)}{\Var(X)}
\]
Since $\Cov(Y,X) = \rho\,\sigma_X\sigma_Y$ and $\Var(X) = \sigma_X^2$:
\[
\frac{\Cov(Y,X)}{\Var(X)} = \frac{\rho\,\sigma_X\sigma_Y}{\sigma_X^2} = \rho\,\frac{\sigma_Y}{\sigma_X}
\]
This is exactly the slope of the conditional mean.  Under joint normality, the BLP coefficient and the CEF slope coincide.

\subsection{The conditional variance is constant}

$\Var(Y\mid X=x) = \sigma_Y^2(1-\rho^2)$ does not depend on $x$.  This is \textbf{homoskedasticity}: the spread of $Y$ around its conditional mean is the same for every value of $X$.

\medskip\noindent
Under joint normality, the key ingredients of the classical linear-normal regression setup follow automatically: a linear CEF, conditionally normal errors, and homoskedasticity.

\subsection{The error variance}

The regression error is $e = Y - \mathbb{E}[Y\mid X]$.  By construction, $\mathbb{E}[e\mid X]=0$.  Because the conditional variance is constant under joint normality:
\[
\Var(e\mid X) = \sigma_Y^2(1-\rho^2)
\]
and therefore
\[
\Var(e) = \mathbb{E}[\Var(e\mid X)] = \sigma_Y^2(1-\rho^2) = \sigma_Y^2 - \frac{[\Cov(Y,X)]^2}{\Var(X)}.
\]
(In general, $\Var(Y\mid X)$ is a random variable measurable with respect to $X$, but here it equals a constant, so $\Var(e)=\Var(e\mid X)$.)

\medskip\noindent
When $|\rho|$ is close to 1, the errors are small and regression explains most of the variation in $Y$.  When $\rho = 0$, the conditional variance equals the unconditional variance $\sigma_Y^2$---knowing $X$ tells us nothing about $Y$.

\section{Does this break down in higher dimensions?}

No.  Everything above generalizes from one regressor to $k$ regressors with the same structure.  The algebra uses block matrix operations instead of scalar fractions, but the logic is identical: partition, condition, complete the square.

\subsection{Setup}

Suppose $(Y, \bm{X})$ are jointly multivariate normal, where $Y$ is a scalar and $\bm{X}$ is a $k\times 1$ vector.  Partition the mean and covariance:
\[
\begin{pmatrix} Y \\ \bm{X} \end{pmatrix}
\sim N\!\left(
\begin{pmatrix} \mu_Y \\ \bm{\mu}_X \end{pmatrix},\
\begin{pmatrix} \sigma_Y^2 & \bm{\Sigma}_{YX} \\ \bm{\Sigma}_{XY} & \bm{\Sigma}_{XX} \end{pmatrix}
\right)
\]
where $\bm{\Sigma}_{YX} = \Cov(Y, \bm{X})$ is $1\times k$ and $\bm{\Sigma}_{XX} = \Var(\bm{X})$ is $k\times k$.

\subsection{The multivariate normal conditional (statement)}

The conditional distribution is:
\begin{tcolorbox}[colback=blue!5, colframe=blue!50]
$Y \mid \bm{X}=\bm{x} \;\sim\; N\!\Big(\mu_{Y|\bm{X}},\;\sigma^2_{Y|\bm{X}}\Big)$

\medskip
\begin{itemize}[leftmargin=2em]
\item \textbf{Conditional mean:}\quad
$\displaystyle\mathbb{E}[Y\mid \bm{X}=\bm{x}] = \mu_Y + \bm{\Sigma}_{YX}\,\bm{\Sigma}_{XX}^{-1}(\bm{x} - \bm{\mu}_X)$

\medskip
\item \textbf{Conditional variance:}\quad
$\displaystyle\Var(Y\mid \bm{X}=\bm{x}) = \sigma_Y^2 - \bm{\Sigma}_{YX}\,\bm{\Sigma}_{XX}^{-1}\,\bm{\Sigma}_{XY}$
\end{itemize}
\end{tcolorbox}

\subsection{Verifying the bivariate case}

When $k=1$, $\bm{\Sigma}_{XX}=\sigma_X^2$, $\bm{\Sigma}_{YX}=\Cov(Y,X)=\rho\,\sigma_Y\sigma_X$, and $\bm{\Sigma}_{XX}^{-1}=1/\sigma_X^2$.  Substituting:
\[
\mathbb{E}[Y\mid X=x] = \mu_Y + \frac{\rho\,\sigma_Y\sigma_X}{\sigma_X^2}(x-\mu_X)
= \mu_Y + \rho\,\frac{\sigma_Y}{\sigma_X}(x-\mu_X) \quad\checkmark
\]
\[
\Var(Y\mid X=x) = \sigma_Y^2 - \frac{(\rho\,\sigma_Y\sigma_X)^2}{\sigma_X^2}
= \sigma_Y^2(1-\rho^2) \quad\checkmark
\]

\subsection{Derivation sketch}

The multivariate derivation follows exactly the same three moves as the bivariate case:

\medskip\noindent
\textbf{Step 1: Form the ratio.}  Write the joint density of $(Y, \bm{X})$ and divide by the marginal density of $\bm{X}$.  Both are multivariate normal densities, so the ratio is an exponential of a difference of quadratic forms:
\[
f_{Y|\bm{X}}(y\mid \bm{x}) \propto
\exp\!\left(
-\frac{1}{2}\left[
\begin{pmatrix} y - \mu_Y \\ \bm{x}-\bm{\mu}_X \end{pmatrix}' \bm{\Sigma}^{-1}
\begin{pmatrix} y - \mu_Y \\ \bm{x}-\bm{\mu}_X \end{pmatrix}
- (\bm{x}-\bm{\mu}_X)'\bm{\Sigma}_{XX}^{-1}(\bm{x}-\bm{\mu}_X)
\right]
\right)
\]
The key tool is the \textbf{block inverse formula}.  For a partitioned positive-definite matrix:
\[
\bm{\Sigma}^{-1} =
\begin{pmatrix} \sigma_Y^2 & \bm{\Sigma}_{YX} \\ \bm{\Sigma}_{XY} & \bm{\Sigma}_{XX} \end{pmatrix}^{-1}
= \begin{pmatrix} \omega^{-1} & -\omega^{-1}\bm{\Sigma}_{YX}\bm{\Sigma}_{XX}^{-1} \\
-\bm{\Sigma}_{XX}^{-1}\bm{\Sigma}_{XY}\omega^{-1} & \bm{\Sigma}_{XX}^{-1}+\bm{\Sigma}_{XX}^{-1}\bm{\Sigma}_{XY}\omega^{-1}\bm{\Sigma}_{YX}\bm{\Sigma}_{XX}^{-1} \end{pmatrix}
\]
where $\omega = \sigma_Y^2 - \bm{\Sigma}_{YX}\bm{\Sigma}_{XX}^{-1}\bm{\Sigma}_{XY}$ is the \emph{Schur complement}---the conditional variance.

\medskip\noindent
\textbf{Step 2: Complete the square.}  Substitute the block inverse and expand.  All terms involving $\bm{x}$ alone (without $y$) cancel between the joint and marginal quadratic forms---exactly as $(x^*)^2$ cancelled in the bivariate case.  What remains is:
\[
-\frac{1}{2\omega}\Big(y - \mu_Y - \bm{\Sigma}_{YX}\bm{\Sigma}_{XX}^{-1}(\bm{x}-\bm{\mu}_X)\Big)^2
\]

\medskip\noindent
\textbf{Step 3: Read off the result.}  This is the kernel of a univariate normal with mean $\mu_Y + \bm{\Sigma}_{YX}\bm{\Sigma}_{XX}^{-1}(\bm{x}-\bm{\mu}_X)$ and variance $\omega$.

\subsection{The same three properties hold}

\begin{enumerate}[leftmargin=2em]
\item \textbf{Linear CEF.}\quad $\mathbb{E}[Y\mid\bm{X}=\bm{x}] = \mu_Y + \bm{\Sigma}_{YX}\bm{\Sigma}_{XX}^{-1}(\bm{x}-\bm{\mu}_X)$ is linear in $\bm{x}$.

\item \textbf{Slope $=$ regression coefficient.}\quad The slope vector is $\bm{\beta} = \bm{\Sigma}_{XX}^{-1}\bm{\Sigma}_{XY} = [\Var(\bm{X})]^{-1}\Cov(\bm{X},Y)$, which is the population multiple regression coefficient.

\item \textbf{Homoskedasticity.}\quad The conditional variance $\sigma_Y^2 - \bm{\Sigma}_{YX}\bm{\Sigma}_{XX}^{-1}\bm{\Sigma}_{XY}$ is a constant---it does not depend on $\bm{x}$.
\end{enumerate}

\medskip\noindent
The bottom line: nothing breaks in higher dimensions.  Joint normality of $(Y, \bm{X})$ implies a linear CEF, normal errors, and homoskedasticity regardless of the number of regressors.  The scalar fractions $\rho\,\sigma_Y/\sigma_X$ and $\sigma_Y^2(1-\rho^2)$ are simply replaced by their matrix counterparts $\bm{\Sigma}_{YX}\bm{\Sigma}_{XX}^{-1}$ and $\sigma_Y^2 - \bm{\Sigma}_{YX}\bm{\Sigma}_{XX}^{-1}\bm{\Sigma}_{XY}$.

\section{But doesn't the normal concentrate on a shell?}

You may have encountered the following fact: if $\bm{Z} = (Z_1, \ldots, Z_k)$ with $Z_j\overset{iid}{\sim}N(0,1)$, then in high dimensions nearly all the probability mass lies near the \emph{surface} of a sphere of radius $\sqrt{k}$, not near the origin.  Specifically:
\[
\|\bm{Z}\|^2 = \sum_{j=1}^k Z_j^2 \sim \chi^2_k, \qquad
\mathbb{E}[\|\bm{Z}\|^2] = k, \qquad
\Var(\|\bm{Z}\|^2) = 2k
\]
As $k$ grows, the standard deviation $\sqrt{2k}$ becomes negligible relative to the mean $k$, so $\|\bm{Z}\|^2/k \to 1$ in probability.  The mass concentrates in a thin shell at radius $\approx\!\sqrt{k}$.

\medskip\noindent
This is real and important in high-dimensional statistics.  But it does \textbf{not} undermine the regression result above, for two reasons.

\subsection{The conditioning derivation is purely algebraic}

The conditional distribution $Y\mid \bm{X}=\bm{x}$ was derived by dividing the joint density by the marginal density.  This is a manipulation of density \emph{functions}---it holds for every value of $\bm{x}$ in $\mathbb{R}^k$, regardless of how probable that particular $\bm{x}$ is.  The formula
\[
\mathbb{E}[Y\mid \bm{X}=\bm{x}] = \mu_Y + \bm{\Sigma}_{YX}\bm{\Sigma}_{XX}^{-1}(\bm{x}-\bm{\mu}_X)
\]
is an identity that holds at every $\bm{x}$, including points near the origin where mass is thin.  Shell concentration tells us \emph{where the data will be}, not that the formula changes at different locations.

\subsection{What concentration \emph{does} affect: finite-sample estimation}

The population result is secure, but \textbf{estimating the regression in finite samples becomes genuinely harder} as $k$ grows, and shell concentration is part of the reason.  This deserves careful treatment, not a hand-wave.

\medskip\noindent
\textbf{The core problem: estimating $\bm{\Sigma}_{XX}^{-1}$.}\quad The population regression slope is $\bm{\beta} = \bm{\Sigma}_{XX}^{-1}\bm{\Sigma}_{XY}$.  OLS estimates this with $\hat{\bm{\beta}} = (\bm{X'X})^{-1}\bm{X'Y}$.  Everything depends on how well the sample covariance matrix $\hat{\bm{\Sigma}}_{XX} = \bm{X'X}/n$ approximates $\bm{\Sigma}_{XX}$.  In low dimensions ($k \ll n$) this is reliable.  In higher dimensions, several things go wrong:

\begin{enumerate}[leftmargin=2em]
\item \textbf{Parameter count outpaces data.}\quad $\bm{\Sigma}_{XX}$ has $k(k+1)/2$ free parameters.  With $k=50$ regressors, that is 1,275 quantities to estimate.  Even with $n=200$ observations, each parameter is pinned down by very little data.

\item \textbf{Eigenvalue spreading (Marchenko--Pastur).}\quad When $k/n$ is not negligible, the eigenvalues of $\hat{\bm{\Sigma}}_{XX}$ spread out relative to those of $\bm{\Sigma}_{XX}$.  The smallest sample eigenvalues are biased toward zero, the largest biased upward.  Since $\hat{\bm{\Sigma}}_{XX}^{-1}$ inverts these eigenvalues, low-variance directions get amplified, injecting noise into $\hat{\bm{\beta}}$.

\item \textbf{Singularity when $k \geq n$.}\quad If $k \geq n$, the sample covariance matrix is literally singular---$\bm{X'X}$ has rank at most $n$, so $(\bm{X'X})^{-1}$ does not exist and OLS cannot be computed at all.
\end{enumerate}

\medskip\noindent
Shell concentration is one geometric manifestation of these high-dimensional difficulties: when $k$ is large, the norms $\|\bm{x}_i - \bm{\mu}_X\|$ concentrate near $\sqrt{k}$, which is a symptom of the same regime where $\hat{\bm{\Sigma}}_{XX}$ becomes unreliable.  But the central econometric issue is the instability of $\hat{\bm{\Sigma}}_{XX}^{-1}$ when $k$ is large relative to $n$, not shell concentration per se.

\medskip\noindent
\textbf{Consequences for OLS.}\quad Even though the population $\bm{\beta}$ is well-defined and the CEF is still linear, the OLS estimator in high dimensions can have:
\begin{itemize}[leftmargin=2em]
\item Enormous variance---individual $\hat{\beta}_j$'s bounce wildly across samples.
\item Poor predictive performance---$\hat{Y} = \bm{x}'\hat{\bm{\beta}}$ overfits the training data.
\item Confidence intervals so wide as to be uninformative, or (worse) nominal coverage that is badly wrong because the $t$- and $F$-distributions assume $k$ is fixed as $n\to\infty$.
\end{itemize}

\medskip\noindent
\textbf{What can be done.}\quad Regularization methods accept some bias in exchange for dramatically lower variance:
\begin{itemize}[leftmargin=2em]
\item \textbf{Ridge regression} shrinks $\hat{\bm{\beta}}$ by replacing $(\bm{X'X})^{-1}$ with $(\bm{X'X}+\lambda\bm{I})^{-1}$, damping the noisy small-eigenvalue directions.
\item \textbf{LASSO} adds an $\ell_1$ penalty that drives some $\hat{\beta}_j$ to exactly zero, performing variable selection.
\item \textbf{Principal component regression} projects $\bm{X}$ onto its leading eigenvectors before regressing, discarding the noisy directions entirely.
\end{itemize}
These methods implicitly acknowledge that the data do not contain enough information to estimate all $k$ slopes reliably.

\subsection{Other effects of high dimensions}

Beyond the estimation of $\bm{\beta}$, shell concentration creates additional difficulties:

\begin{itemize}[leftmargin=2em]
\item \textbf{Distance-based methods.}  In high dimensions, pairwise distances become less informative, so methods that rely on finding ``close'' neighbors (e.g., kernel regression, $k$-NN) often deteriorate.

\item \textbf{Density evaluation.}  The joint density $f(\bm{x})$ evaluated at any particular point goes to zero as $k\to\infty$, even on the shell.  This makes likelihood-based model comparison in very high dimensions difficult.  The \emph{conditional} density for the scalar $Y$ given $\bm{X}=\bm{x}$ remains a well-behaved univariate normal, but evaluating the marginal likelihood requires integrating over the high-dimensional shell.
\end{itemize}

\subsection{Bottom line}

The population-level result is airtight: under joint normality, the CEF is linear with normal homoskedastic errors for any $k$.  But the \emph{estimation} problem degrades seriously as $k$ grows relative to $n$, because the sample covariance matrix becomes a noisy estimate of $\bm{\Sigma}_{XX}$ and inverting it amplifies that noise.  Shell concentration is one geometric symptom of this high-dimensional regime, not the root cause.  The distinction matters---the \emph{model} is fine, but our ability to \emph{learn the model from data} is not, and recognizing this gap is the starting point for regularization and high-dimensional inference.

\section{Summary}

Starting from the bivariate normal density, we showed:
\begin{enumerate}[leftmargin=2em]
\item The conditional distribution $Y\mid X=x$ is normal (Section~4).
\item The conditional mean is linear in $x$, with slope $\rho\,\sigma_Y/\sigma_X = \Cov(Y,X)/\Var(X)$ (Section~5.1--5.2).
\item The conditional variance $\sigma_Y^2(1-\rho^2)$ is constant in $x$ (Section~5.3).
\end{enumerate}
Under joint normality, the CEF \emph{is} the population linear regression, and the classical normal regression assumptions hold by construction.  Section~6 shows that the same result extends to $k$ regressors via the partitioned multivariate normal, with scalar operations replaced by their matrix analogues.  Section~7 explains why the concentration-on-a-shell phenomenon in high dimensions does not affect these results: the conditional distribution is derived pointwise for any $\bm{x}$, and the practical difficulties of high dimensions are estimation problems, not failures of the underlying population relationship.

\end{document}
